\section{Problem Formulations}\label{chap:problem}

\subsection{The Primal Problem (P)} % TODO
We begin with the first soft-margin classification problem we are interested in solving, which is
formulated as a convex optimization problem. This will serve as the primal problem in our analysis.

\begin{align*}\label{eq:primal-problem}
	\min_{\symbf{w}, b, \symbf{\xi}} \quad
	 & f(\symbf{w}, b, \symbf{\xi})
	= \frac{1}{2}\|\symbf{w}\|_2^2 + C\,\symbf{1}_M^\top \symbf{\xi}
	\quad
	\text{subject to}
	\quad
	\begin{cases}
		\symbf{1}_M - Y\bigl(X\,\symbf{w} + b\,\symbf{1}_M\bigr) - \symbf{\xi} & \le \symbf{0}_M, \\[6pt]
		\symbf{\xi} \;\;                                                       & \ge \symbf{0}_M,
	\end{cases}\tag{P}
\end{align*}

Here, $\symbf{w} \in \R^d$ represents the weight vector, $b \in \R$ is the bias term, and $\symbf{\xi} \in \R^M$ is the vector of slack variables.
The data matrix is $X = [\symbf{x}_1, \ldots, \symbf{x}_M] \in \R^{d \times M}$ with labels $\symbf{y}=[y_1, \ldots, y_M]^\top \in \{\pm1\}^M$ and diagonal label matrix $Y = \operatorname{diag}(y_1,\ldots,y_M)$.
The parameter $C>0$ controls the penalty for misclassification.

\subsubsection{Existence and Uniqueness of Constrained Optimization Problems}
To prove the existence and uniqueness of a minimizer for a constrained optimization problem,
three properties are required. Feasibility, Coercivity, and Convexity. Feasibility ensures that the constraint set is non-empty, meaning
that there is at least one point that satisfies all constraints. Coercivity of the objective
function guarantees that the function tends to infinity as the norm of the variable tends to infinity,
ensuring the existence of a global minimum on the feasible set. Finally, convexity of the objective
function implies that any local minimizer is also a global one, and if the function is strictly convex,
it ensures that the minimizer is unique. Therefore, the combination of feasibility, coercivity, and
strict convexity, establishes both the existence and uniqueness of a global minimizer.

\begin{theorem}{Existence \& Uniqueness of Primal Solution}{existence-uniqueness-primal}
  The primal problem \eqref{eq:primal-problem} always admits at least one global minimizer. Furthermore, the weight vector $\symbf{w}^\star$ is unique due to strict convexity in $\symbf{w}$. However, $b^\star$ and/or $\symbf{\xi}^\star$ may fail to be unique under certain degenerate conditions.
\end{theorem}

\begin{proof}[Proof of Theorem~\ref{thm:existence-uniqueness-primal}]
  \paragraph{Existence of a solution.}
  The primal problem~\eqref{eq:primal-problem} is a convex quadratic program with linear constraints. We check two key facts:
  \begin{itemize}
    \item \textbf{Feasibility:} The point
    \[
      \symbf{w} = \symbf{0}_d,\quad
      b = 0,\quad
      \symbf{\xi} = \symbf{1}_M
    \]
    satisfies
    \(
      \symbf{1}_M - Y\bigl(X\,\symbf{0}_d + 0\cdot\symbf{1}_M\bigr) - \symbf{1}_M = \symbf{0}_M \le \symbf{0}_M,
      \quad
      \symbf{1}_M \ge \symbf{0}_M.
    \)
    Hence the feasible set is nonempty.
    \item \textbf{Coercivity:} The objective
    \[
       f(\symbf{w},b,\symbf{\xi}) 
       = \tfrac12\,\|\symbf{w}\|_2^2 
         \;+\; C\,\bigl(\symbf{1}_M^\top \symbf{\xi}\bigr)
    \]
    grows unboundedly whenever $\|\symbf{w}\|\to\infty$ or $\|\symbf{\xi}\|\to\infty$ ($C>0$). So the objective cannot go to $-\infty.$ By standard convexity arguments, a global minimizer must exist.
  \end{itemize}

  \paragraph{Uniqueness of a solution.}
  \begin{itemize}
    \item \textbf{Uniqueness of $\symbf{w}^\star$:}
    The strictly convex term $\tfrac12\,\|\symbf{w}\|_2^2$ ensures that $\symbf{w}^\star$ is unique once the problem is solved, because no other $\symbf{w}$ can yield the same minimal value unless it coincides with $\symbf{w}^\star$.
    \item \textbf{Possible non-uniqueness of $b^\star$:}
    Since $b$ does not appear in the quadratic term, it may not be pinned down uniquely by the constraints (unless certain support vectors fix it via KKT conditions). In symmetric or degenerate configurations, $b$ can shift while adjusting $\symbf{w}$ accordingly.
    \item \textbf{Possible non-uniqueness of $\symbf{\xi}^\star$:}
    For points with $\alpha_i^\star=0$, one only needs $\xi_i \ge 0$, so there may be infinitely many ways to satisfy that constraint. Conversely, for $\alpha_i^\star = C$ or $0<\alpha_i^\star<C$, $\xi_i^\star$ is determined exactly by the constraints. In degenerate cases, overall $\symbf{\xi}^\star$ need not be unique.
  \end{itemize}
\end{proof}
\begin{theorem}{Existence \& Uniqueness of Primal Solution}{existence-uniqueness-primal}
	The primal problem \eqref{eq:primal-problem} always admits at least one global minimizer. Moreover, the weight vector $\symbf{w}^\star$ is unique due to strict convexity of the objective in $\symbf{w}$. However, the bias $b^\star$ and slack variables $\symbf{\xi}^\star$ may fail to be unique under certain degenerate configurations.
\end{theorem}

\begin{proof}[Proof of Theorem~\ref{thm:existence-uniqueness-primal}]
	\paragraph{Existence of a solution.}
	The primal problem~\eqref{eq:primal-problem} is a convex quadratic program with affine inequality constraints. To establish existence of a minimizer, we verify feasibility and coercivity of the objective.

	\begin{itemize}
		\item \textbf{Feasibility:} Let $\symbf{w} = \symbf{0}_d$, $b = 0$, and $\symbf{\xi} = \symbf{1}_M$. Then,
		      \[
			      \symbf{1}_M - Y\left(X\symbf{w} + b\,\symbf{1}_M\right) - \symbf{\xi}
			      = \symbf{1}_M - Y(\symbf{0}_M) - \symbf{1}_M = \symbf{0}_M \le \symbf{0}_M,
		      \]
		      and clearly $\symbf{\xi} = \symbf{1}_M \ge \symbf{0}_M$. Thus, this choice satisfies all constraints, and the feasible region is non-empty.

		\item \textbf{Coercivity:} The objective function
		      \[
			      f(\symbf{w}, b, \symbf{\xi}) = \frac{1}{2} \|\symbf{w}\|_2^2 + C\, \symbf{1}_M^\top \symbf{\xi}
		      \]
		      is coercive in $\symbf{w}$ and $\symbf{\xi}$. As either $\|\symbf{w}\| \to \infty$ or $\|\symbf{\xi}\| \to \infty$, the objective tends to infinity. Since the feasible region is non-empty and convex, and the objective is continuous and coercive, the problem attains a minimum.
	\end{itemize}

	\paragraph{Uniqueness of a solution.}
	\begin{itemize}
		\item \textbf{Uniqueness of $\symbf{w}^\star$:}
		      The term $\tfrac{1}{2}\|\symbf{w}\|^2$ is strictly convex in $\symbf{w}$, while the remaining terms are affine. As a result, the full objective is strictly convex in $\symbf{w}$, guaranteeing that $\symbf{w}^\star$ is unique.

		\item \textbf{Possible non-uniqueness of $b^\star$:}
		      The variable $b$ appears only linearly in the constraints and not in the objective. In cases where the constraints do not uniquely determine $b$—for example, when multiple values of $b$ yield the same classification margin—a degree of freedom in $b$ remains. This can happen, for instance, when there are no support vectors exactly on the margin.

		\item \textbf{Possible non-uniqueness of $\symbf{\xi}^\star$:}
		      The slack variables $\xi_i$ only appear linearly in both the objective and the constraints. For points that are not support vectors (i.e., not active in the constraint), the optimal $\xi_i^\star$ may be any non-negative value satisfying the constraint, leading to potential non-uniqueness. However, for support vectors that lie within or on the margin, $\xi_i^\star$ becomes uniquely determined.
	\end{itemize}
\end{proof}
